% ── Dokumenty ─────────────────────────────────────────────────
\chapter{Dokumenty}

        Správa dokumentov je jednou z~kľúčových funkcionalít systému. Umožňuje
nahrávanie, organizovanie a~sprístupňovanie súborov priamo v~kontexte projektov
a~úloh. Systém podporuje formáty PDF a~DOCX, pričom z~každého nahraného súboru
automaticky extrahuje text pre potreby AI~chatu.

\section{Tvorba dokumentov}

        Nahrávanie dokumentu je dostupné z~detailu projektu alebo úlohy. Používateľ
vyberie súbor zo svojho zariadenia a~systém vykoná nasledujúce kroky:

\begin{itemize}
  \item \textbf{Validácia} -- systém overí, či má súbor povolený formát (PDF alebo DOCX)
        a~či neprekračuje maximálnu veľkosť 16~MB.
  \item \textbf{Uloženie} -- súbor je uložený na lokálny filesystém servera do adresára
        \texttt{uploads/} s~unikátnym názvom (UUID), aby sa zabránilo kolíziám mien.
  \item \textbf{Extrakcia textu} -- z~PDF súborov je text extrahovaný knižnicou
        PyPDF2, z~DOCX súborov knižnicou python-docx. Extrahovaný text je uložený
        do stĺpca \texttt{extracted\_text} v~databáze.
  \item \textbf{Záznam v~databáze} -- vytvorí sa záznam \texttt{Document} s~metadátami:
        pôvodný názov súboru, cesta k~uloženému súboru, typ, projekt a~nahrávateľ.
\end{itemize}

\begin{figure}[H]
  \centering
  % \includegraphics[width=0.85\textwidth]{images/nahravanie_dokumentu.png}
  \fbox{\parbox{0.80\textwidth}{\centering\vspace{2.5cm}
    [TBD]
  \vspace{2.5cm}}}
  \caption{Nahrávanie dokumentu}
  \label{fig:nahravanie}
\end{figure}

        Nahrané dokumenty sú zobrazené v~prehľadnom zozname s~názvom, typom súboru
a~dátumom nahratia. Každý dokument je možné stiahnuť alebo otvoriť v~AI~chate
na kladenie otázok nad jeho obsahom.

\begin{figure}[H]
  \centering
  % \includegraphics[width=0.9\textwidth]{images/zoznam_dokumentov.png}
  \fbox{\parbox{0.85\textwidth}{\centering\vspace{2.5cm}
    [TBD]
  \vspace{2.5cm}}}
  \caption{Zoznam dokumentov projektu}
  \label{fig:dokumenty}
\end{figure}

\section{Zdieľanie dokumentov}

        Dokumenty sú zdieľané v~rámci projektu -- všetci členovia projektu majú
prístup k~dokumentom nahraným v~jeho kontexte. Tým je zabezpečená tímová
spolupráca bez nutnosti posielania súborov e-mailom alebo cez externé úložiská.

        Systém neposiela obsah dokumentov na externé servery. Súbory sú uložené
lokálne na serveri organizácie a~AI~model pristupuje k~extrahovanému textu
priamo z~lokálnej databázy. Toto riešenie zaručuje, že citlivé informácie zostanú výhradne
v infraštruktúre organizácie.
